\todo{Lecture 21}
\chapter{Systèmes différentielles linéaires}
On veut résoudre un système
\begin{align*}
x_{1}'(t) & =a_{11}x_{1}(t)+\dots +a_{1n}x_{n}(t) \\
 & ~~\vdots \\
x'_{n}(t)  & =a_{n_{1}}x_{1}(t)+\dots +a_{nn}x_{n}(t)
\end{align*}
ou $A=(a_{ij})^{n}_{i,j=1}\in \mathbf{R}^{n\times n}$. On cherche $x_{i}:\mathbf{R}\to \mathbf{R},t\mapsto x_{i}(t)$ derivable telle que le système est satisfait. Etant données des conditions initiales 
$$
x(0)=\begin{pmatrix}
x_{1}(0) \\
\vdots \\
x_{n}(0)
\end{pmatrix}
$$
on veut que 
$$
x(t)=\begin{pmatrix}
x_{1}(t) \\
\vdots \\
x_{n}(t)
\end{pmatrix}
$$
respecte ces condition initiales.
\begin{remark}
On sait (c.f. cours d'analyse) qu'il existe une telle solution $x(t)$ du système telle que les conditions initiales sont satisfaites.
\end{remark}
\begin{example}\label{exe2101}
Soit
$$
A=\begin{pmatrix}
0 & -1  \\
1 & 0
\end{pmatrix},
$$
on cherche
$$
x(t)=\begin{pmatrix}
x_{1}(t) \\
x_{2}(t)
\end{pmatrix}
$$
telle que 
$$
\begin{pmatrix}
x_{1}'(t) \\
x_{2}'(t)
\end{pmatrix}=Ax(t)=\begin{pmatrix}
-x_{2}(t) \\
x_{1}(t)
\end{pmatrix}.
$$
On a une solution
$$
\hat{x}_{1}(t)=\cos t(t),\quad \hat{x}_{2}(t)=\sin (t)
$$
et une autre solution
$$
\overline{x}_{1}(t)=-\sin (t),\quad \overline{x}_{2}(t)=\cos(t).
$$
On observe que ces deux solution sont différentes ($\hat{x}(t)\neq\alpha \overline{x}(t),~ \alpha \in \mathbf{R}$) et que pour tous $\alpha,\beta \in \mathbf{R}$ leur combinaison linéaire $\alpha \hat{x}(t)+\beta \,\overline{x}(t)$ est aussi une solution. 

On suppose les conditions initiales $x(0)=(1/2,5)$, alors 
$$
x(t)=\frac{1}{2}\hat{x}(t)+5\, \overline{x}(t)
$$
est une solution du système respectant les conditions initiales.
\end{example}
\begin{reminder}
On a l'exponentiel 
\begin{align*}
e^{x}:\mathbf{C} & \to \mathbf{C} \\
x & \mapsto \sum_{i=0}^{\infty} \frac{x^{i}}{i!} 
\end{align*}
qui est holomorphe et $(e^{x})'=e^{x}$.

Pour $x=a+ib\in \mathbf{C}$. On a 
$$
e^{x}=e^{ a+ib }=e^{a}(e^{ ib })=e^{ a }(\cos(b)+i\sin(b)).
$$
\end{reminder}

Soit $A\in \mathbf{C}^{n\times n}$. On a
\begin{align*}
e^{A} & =I + \frac{A}{1!}+\frac{A^{2}}{2!}+\dots \\
 & =\sum_{i=0}^{\infty} \frac{A^{i}}{i!}
\end{align*}
\begin{question}
Est-ce que cette série converge ?
\end{question}
\begin{answer}
On a que cette série converge si pour tout $\varepsilon>0$, il existe $N\in \mathbf{N}$ tel que $\forall m\ge N$ :
$$
b(m):= \sum_{i=0}^{m} \frac{A^{i}}{i!}\in \mathrm{Voisinage}_{\varepsilon}(e^{A})
$$
ou le voisinage est relatif a $\lVert \cdot \rVert_{F}\sim \lVert \cdot \rVert_{2}$ en $\mathbf{R}^{n^{2}}$. On veut montrer que $b(m)$ est une suite de Cauchy c-a-d
$$
\forall\varepsilon>0\,\exists N\in \mathbf{N}\,\forall m>n\ge N: \underbrace{ \left\lVert  \sum_{i=n+1}^{m} \frac{A^{i}}{i!}  \right\rVert _{F}}_{ =\left\lVert  b(m)-b(n)\right\rVert _{F} }<\varepsilon.
$$
On a 
$$
\left\lVert  \sum_{i=n+1}^{m} \frac{A^{i}}{i!}  \right\rVert _{F} \le \sum_{i=n+1}^{m} \frac{\lVert A^{i} \rVert _{F}}{i!}
$$
\begin{lemma}
Pour $A,B\in \mathbf{C}^{n\times n}$ on a 
$$
\lVert AB \rVert _{F}\le\lVert A \rVert_{F}\lVert B\rVert _{F}.
$$
\end{lemma}
\begin{proof}
Soient
$$
A=\begin{pmatrix}
a_{1}^{\mathsf{T}} \\
\vdots \\
a_{n}^{\mathsf{T}}
\end{pmatrix},\quad B=(\overline{b}_{1},\dots,\overline{b}_{n}).
$$
On a $(a_{1}^{\mathsf{T}}\overline{b}_{1})(\overline{a}_{1}^{\mathsf{T}}b_{1})=|a_{1}^{\mathsf{T}}\overline{b}_{1}|^{2}$. Donc
\begin{align*}
\lVert AB \rVert _{F}^{2} & =\sum _{i,j=1}^{n}(a_{i}^{\mathsf{T}}\overline{b}_{j})(\overline{a}_{i}^{\mathsf{T}}b_{j}) \\
 & \le \sum_{i,j=1}^{n} \lVert a_{i} \rVert ^{2}\lVert b_{j} \rVert ^{2} \\
 & =\left( \sum_{i=1}^{n} \lVert a_{i} \rVert ^{2} \right) \left( \sum_{j=1}^{n} \lVert b_{i} \rVert ^{2} \right) \\
 & =\lVert A \rVert^{2} _{F}\lVert B \rVert ^{2}_{F}.
\end{align*}
\end{proof}

On sait que 
$$
c(n)=\sum_{i=0}^{n} \frac{\lVert A \rVert _{F}^{i}}{i!}
$$
est de Cauchy dans le sens de la suite $c:\mathbf{N}_{0}\to \mathbf{R}_{\ge 0}$. Alors
$$
\forall\varepsilon>0\,\exists N_{\varepsilon}\in \mathbf{N}\,\forall m,n\ge N_{\varepsilon}:|c(n)-c(m)|<\varepsilon.
$$
Avec cet $N_{\varepsilon}$ on a
$$
\left\lVert  \sum_{i=n+1}^{m} \frac{A^{i}}{i!}  \right\rVert _{F}<\varepsilon
$$
donc la série intégrale converge.
\end{answer}
\begin{remark}
Soit $A\in \mathbf{C}^{m\times n}$. La norme Frobenius de $A$ est définie comme
$$
\lVert A \rVert _{F}=\left( \sum _{i,j=1}^{n}a_{ij}\overline{a}_{ij} \right)^{1/2}
$$
qui est équivalente a somme de normes euclidiennes des colonnes de $A$.
\end{remark}

Soit $A\in \mathbf{R}^{n\times n}$ une matrice. On a pour tout $t\in \mathbf{R}$ :
$$
e^{ tA }=\sum_{i=0}^{\infty} t^{i} \frac{A^{i}}{i!}.
$$
Toutes composantes de $e^{ tA }$ sont des fonctions en $t$ et sont representes comme une serie integrale
$$
f(t)=\sum_{i=0}^{\infty} t^{i} \frac{(A^{i})_{kj}}{i!} 
$$
avec derivee
$$
f'(t)=\sum_{i=0}^{\infty} t^{i}  \frac{(A(A^{i}))_{ij}}{i!}.  
$$
Toutes fonctions dans $e^{ tA }$ dérivée par rapport a $t$
$$
\frac{d}{dt}(e^{ tA })=Ae^{ tA }.
$$
\begin{theorem}
La solution d'un système différentiel linéaire qui respecte les conditions initiales $x(0)\in \mathbf{R}^{n}$ est 
$$
x(t)=e^{ tA }x(0).
$$
\end{theorem}
\begin{proof}
On a 
$$
e^{0A}x_{0}=I_{n}x_{0}
$$
et 
$$
x'(t)=Ae^{ tA }x(0)=Ax(t).
$$
\end{proof}
\begin{example}[Continuation du exemple \ref{exe2101}]
On calcule
$$
e^{ tA }=\sum_{i=0}^{\infty}t^{i} \frac{A^{i}}{i!}. 
$$
On a 
$$
A^{0}=I, A^{1}=A, A^{2}=-I,A^{3}=-A,A^{4}=I,\dots
$$
donc
$$
A^{2k}=(-1)^{k}I,\quad A^{2k+1}=\begin{pmatrix}
0 & (-1)^{k+1} \\
(-1)^{k} & 0
\end{pmatrix}.
$$
Alors on a 
$$
e^{ tA }=\begin{pmatrix}
\sum_{k=0}^{\infty} \frac{(-1)^{k}t^{2k}}{(2k)!} & \sum_{k=0}^{\infty} \frac{(-1)^{k+1}t^{2k+1}}{(2k+1)!}  \\
\sum_{k=0}^{\infty} \frac{(-1)^{k}t^{2k+1}}{(2k+1)!} & \sum_{k=0}^{\infty}  \frac{(-1)^{k}t^{2k}}{(2k)!}
\end{pmatrix}
$$
ou on reconnait les power series de $\cos(t)$ et $\sin(t)$, donc
$$
e^{ tA }=\begin{pmatrix}
\cos(t) & -\sin(t) \\
\sin(t) & \cos(t)
\end{pmatrix}
$$
qui verifie
$$
e^{ tA }\begin{pmatrix}
\frac{1}{2} \\
5
\end{pmatrix}=\frac{1}{2}\begin{pmatrix}
\cos(t) \\
\sin(t)
\end{pmatrix}+5\begin{pmatrix}
-\sin (t) \\
\cos(t)
\end{pmatrix}.
$$
\end{example}
